\section{Introduction}
\label{sec:introduction}

Statistical power is the probability that a test correctly rejects a
false null hypothesis.
In redistricting, the null hypothesis is that the enacted plan is a
random draw from the space of valid plans (i.e., it was not gerrymandered).
The alternative is that the plan was selected to produce a specific
partisan advantage.

Power depends on three quantities:
(1) the significance level $\alpha$ (typically 0.05 for scientific
claims, or the court's threshold for ``clear and convincing'' evidence);
(2) the effect size $\delta$ (the partisan advantage in percentage points
that the test is designed to detect);
and (3) the effective sample size $\ess$ of the ensemble
(which determines how precisely the null distribution can be estimated).

Existing redistricting analyses report none of these quantities.
They run a chain of $n$ steps, compute the enacted plan's percentile,
and report it.
The implicit power guarantee is undefined: a chain that has not converged
may have essentially zero power to detect any effect, regardless of $n$.

This paper derives the power function explicitly and translates it into
minimum ensemble sizes for each state in the 2020 congressional redistricting.

\subsection{Why Power Matters in Court}

A defendant in a redistricting challenge can argue:
``The expert's ensemble was too small to reliably detect the alleged
gerrymander.
Even if the enacted plan is at the 1st percentile, an underpowered
ensemble cannot distinguish between a genuine gerrymander and random
variation.''

This argument is technically sound if the ensemble is indeed underpowered.
A power-analysis defense requires the plaintiff's expert to demonstrate
that the ensemble has adequate power to detect the alleged effect before
the court can rely on the percentile claim.

Conversely, an expert who demonstrates 80\% power at the observed effect
size provides a defensible claim: ``Had this plan been drawn randomly,
we would have 80\% probability of detecting it as an outlier at $\alpha = 0.05$.
We did detect it. The probability of a false positive is at most 5\%.''

\subsection{Road Map}

Section~\ref{sec:power} derives the power function from first principles.
Section~\ref{sec:ess} translates the ESS requirement into chain step requirements
using G.4's autocorrelation estimates.
Section~\ref{sec:table} provides the 50-state minimum ensemble table.
Section~\ref{sec:current} evaluates current practice against the new standards.
Section~\ref{sec:litigation} gives the litigation-grade recommendations.
Section~\ref{sec:conclusion} concludes.
